\documentclass[11pt,a4paper]{ctexart}
\usepackage[margin=2.2cm]{geometry}
\usepackage{xcolor}
\usepackage{enumitem}
\usepackage{amssymb}
\usepackage{titlesec}
\usepackage{hyperref}
\definecolor{navy}{HTML}{1A365D}
\definecolor{ink}{HTML}{1F2933}
\definecolor{muted}{HTML}{66788A}
\definecolor{soft}{HTML}{EEF4FA}
\definecolor{gold}{HTML}{8A5A12}
\hypersetup{colorlinks=true,linkcolor=navy,urlcolor=navy}
\setlist{nosep}
\titleformat{\section}{\Large\bfseries\color{navy}}{}{0pt}{}
\titleformat{\subsection}{\large\bfseries\color{navy}}{}{0pt}{}
\pagestyle{plain}

\title{\textbf{宾语从句专项}\\[4pt]
\large 人教版（新课纲版）九年级上册 · Unit 3 Smart Learning / Unit 4 Our Memory}
\author{初中英语资料站}
\date{2026-11-07}

\begin{document}
\maketitle

\section*{使用说明}
\begin{itemize}[leftmargin=1.4em]
\item 本讲义对应人教版（新课纲版）九年级上册 Unit 3、Unit 4 的宾语从句要求：引导词、陈述语序、时态呼应。
\item 共 6 讲语法、6 级 48 题；建议按“基础 → 语序 → 时态 → U3 → U4 → 中考综合”的顺序完成。
\item 在线版：\url{https://1901151458-prog.github.io/english-materials/object-clauses-g9/}
\end{itemize}

\section{宾语从句是什么}
宾语从句是用一个完整句子充当动词的宾语。它有三要素：\textbf{引导词、语序、时态}。
\begin{itemize}[leftmargin=1.4em]
\item I know \underline{that he is right}.（that he is right 作 know 的宾语）
\item I wonder \underline{if he will come}.（if he will come 作 wonder 的宾语）
\item Do you know \underline{where he lives}?（where he lives 作 know 的宾语）
\end{itemize}

\section{三类引导词}
\begin{center}
\begin{tabular}{p{3.2cm}p{3.2cm}p{7.4cm}}
\hline
\textbf{原句类型} & \textbf{引导词} & \textbf{说明}\\
\hline
陈述句 & that & 口语中可省略；从句意思完整。\\
一般疑问句 & if / whether & 表示“是否”；or not 连用、介词后、不定式前多用 whether。\\
特殊疑问句 & wh- 词 & 用原疑问词，但后面必须改为陈述语序。\\
选择疑问句 & whether ... or ... & 常与 or not 搭配。\\
\hline
\end{tabular}
\end{center}

\section{陈述语序}
宾语从句一律用“主语 + 谓语”的陈述语序，不能保留疑问句的倒装。
\begin{itemize}[leftmargin=1.4em]
\item \textcolor{red}{×} Do you know where does he live?
\item \textcolor{green!50!black}{✓} Do you know where he lives?
\item \textcolor{red}{×} I wonder what is he doing.
\item \textcolor{green!50!black}{✓} I wonder what he is doing.
\end{itemize}

\section{时态呼应}
\begin{center}
\begin{tabular}{p{3.6cm}p{4.4cm}p{7.2cm}}
\hline
\textbf{主句时态} & \textbf{从句时态} & \textbf{例句}\\
\hline
现在时 / 将来时 & 按实际需要 & She says (that) she is busy.\\
过去时 & 一般过去时 & He said (that) he was busy.\\
过去时 & 过去进行时 & She asked what I was doing.\\
过去时 & 过去完成时 & He said he had finished it.\\
过去时 & 过去将来时 & I didn't know when he would come.\\
过去时 & 客观真理仍用一般现在时 & The teacher told us that light travels faster than sound.\\
\hline
\end{tabular}
\end{center}
口诀：\textbf{主现从随便，主过从过去，真理用现在。}

\section{五个高频易错点}
\begin{enumerate}[leftmargin=1.4em]
\item 直接把疑问句语序搬进从句：× where does he live → ✓ where he lives。
\item 主句过去时，从句还用现在时：× He said he is busy. → ✓ He said he was busy.（客观真理除外）
\item whether / if 混用：× if ... or not → ✓ whether ... or not。
\item 把 if 条件句当成宾语从句：I don't know if he will come.（是否） / If he comes, I will tell you.（如果）
\item 特殊疑问词后多写主语：× who did he help → ✓ who he helped。
\end{enumerate}

\section{对应 Unit 3 / Unit 4}
\begin{itemize}[leftmargin=1.4em]
\item \textbf{Unit 3 Smart Learning}：I wonder how I can improve my reading speed. / Our teacher told us that good learners often ask questions. / Do you know if the library has a quiet study room?
\item \textbf{Unit 4 Our Memory}：I remember that I left my keys on the desk. / Do you know where I put my phone? / She asked me whether I could remember her name.
\end{itemize}

\newpage
\section{分级练习（48 题）}
\noindent 按单元与难度排列，先做基础与语序，再做时态与教材语境，最后做中考综合题。

\subsection*{第 1 级 · 基础 · 选引导词}
\noindent\textcolor{gray}{陈述句、一般疑问句、特殊疑问句分别用什么引导词。}\par\smallskip
\begin{enumerate}[leftmargin=*,itemsep=4pt]
\item I think \_\_\_ English is interesting.
\begin{enumerate}[label=\Alph*.,leftmargin=2em,itemsep=1pt,topsep=2pt]
\item that
\item if
\item what
\item how
\end{enumerate}
\item She asked \_\_\_ I could help her.
\begin{enumerate}[label=\Alph*.,leftmargin=2em,itemsep=1pt,topsep=2pt]
\item that
\item if
\item what
\item who
\end{enumerate}
\item Can you tell me \_\_\_ the library is?
\begin{enumerate}[label=\Alph*.,leftmargin=2em,itemsep=1pt,topsep=2pt]
\item that
\item if
\item where
\item whether
\end{enumerate}
\item I don't know \_\_\_ he will come or not.
\begin{enumerate}[label=\Alph*.,leftmargin=2em,itemsep=1pt,topsep=2pt]
\item that
\item if
\item whether
\item what
\end{enumerate}
\item He told me \_\_\_ he had finished the work.
\begin{enumerate}[label=\Alph*.,leftmargin=2em,itemsep=1pt,topsep=2pt]
\item that
\item if
\item what
\item how
\end{enumerate}
\item Do you know \_\_\_ the meeting will start?
\begin{enumerate}[label=\Alph*.,leftmargin=2em,itemsep=1pt,topsep=2pt]
\item that
\item if
\item when
\item whether
\end{enumerate}
\item I wonder \_\_\_ we should do next.
\begin{enumerate}[label=\Alph*.,leftmargin=2em,itemsep=1pt,topsep=2pt]
\item that
\item if
\item what
\item how
\end{enumerate}
\item Could you tell me \_\_\_ you like the book?
\begin{enumerate}[label=\Alph*.,leftmargin=2em,itemsep=1pt,topsep=2pt]
\item that
\item if
\item why
\item what
\end{enumerate}
\end{enumerate}

\subsection*{第 2 级 · 语序 · 陈述语序}
\noindent\textcolor{gray}{宾语从句一律用“主语 + 谓语”，不能倒装。}\par\smallskip
\begin{enumerate}[leftmargin=*,itemsep=4pt]
\item Do you know \_\_\_?
\begin{enumerate}[label=\Alph*.,leftmargin=2em,itemsep=1pt,topsep=2pt]
\item where he lives
\item where does he live
\item where he live
\item where is he live
\end{enumerate}
\item Could you tell me \_\_\_?
\begin{enumerate}[label=\Alph*.,leftmargin=2em,itemsep=1pt,topsep=2pt]
\item when does the train leave
\item when the train leaves
\item when the train leave
\item when is the train leave
\end{enumerate}
\item I don't know \_\_\_.
\begin{enumerate}[label=\Alph*.,leftmargin=2em,itemsep=1pt,topsep=2pt]
\item what is he doing
\item what he is doing
\item what does he do
\item what he doing
\end{enumerate}
\item She asked me \_\_\_.
\begin{enumerate}[label=\Alph*.,leftmargin=2em,itemsep=1pt,topsep=2pt]
\item if I was okay
\item if was I okay
\item that I was okay
\item if I am okay
\end{enumerate}
\item Can you tell me \_\_\_?
\begin{enumerate}[label=\Alph*.,leftmargin=2em,itemsep=1pt,topsep=2pt]
\item how can I get there
\item how I can get there
\item how could I get there
\item how I could get there
\end{enumerate}
\item I wonder \_\_\_.
\begin{enumerate}[label=\Alph*.,leftmargin=2em,itemsep=1pt,topsep=2pt]
\item why was she late
\item why she was late
\item why did she late
\item why she late
\end{enumerate}
\item Do you know \_\_\_?
\begin{enumerate}[label=\Alph*.,leftmargin=2em,itemsep=1pt,topsep=2pt]
\item who wrote this book
\item who did write this book
\item who this book wrote
\item who was wrote this book
\end{enumerate}
\item He asked me \_\_\_.
\begin{enumerate}[label=\Alph*.,leftmargin=2em,itemsep=1pt,topsep=2pt]
\item whether I had seen the film
\item whether had I seen the film
\item that I had seen the film
\item if I have seen the film
\end{enumerate}
\end{enumerate}

\subsection*{第 3 级 · 时态 · 呼应规则}
\noindent\textcolor{gray}{主句时态与从句时态如何呼应，客观真理除外。}\par\smallskip
\begin{enumerate}[leftmargin=*,itemsep=4pt]
\item He said (that) he \_\_\_ busy then.
\begin{enumerate}[label=\Alph*.,leftmargin=2em,itemsep=1pt,topsep=2pt]
\item is
\item was
\item will be
\item has been
\end{enumerate}
\item She asks me if I \_\_\_ ready.
\begin{enumerate}[label=\Alph*.,leftmargin=2em,itemsep=1pt,topsep=2pt]
\item am
\item was
\item will
\item had
\end{enumerate}
\item The teacher told us that light \_\_\_ faster than sound.
\begin{enumerate}[label=\Alph*.,leftmargin=2em,itemsep=1pt,topsep=2pt]
\item travels
\item traveled
\item will travel
\item has traveled
\end{enumerate}
\item I didn't know whether he \_\_\_ come.
\begin{enumerate}[label=\Alph*.,leftmargin=2em,itemsep=1pt,topsep=2pt]
\item will
\item would
\item is
\item has
\end{enumerate}
\item She wanted to know what I \_\_\_ doing.
\begin{enumerate}[label=\Alph*.,leftmargin=2em,itemsep=1pt,topsep=2pt]
\item am
\item was
\item will be
\item have
\end{enumerate}
\item He asked if I \_\_\_ ever been to Beijing.
\begin{enumerate}[label=\Alph*.,leftmargin=2em,itemsep=1pt,topsep=2pt]
\item have
\item had
\item was
\item am
\end{enumerate}
\item “Where is he?” She asked where he \_\_\_.
\begin{enumerate}[label=\Alph*.,leftmargin=2em,itemsep=1pt,topsep=2pt]
\item is
\item was
\item has been
\item will be
\end{enumerate}
\item “Do you like music?” He asked me \_\_\_.
\begin{enumerate}[label=\Alph*.,leftmargin=2em,itemsep=1pt,topsep=2pt]
\item if I like music
\item if I liked music
\item that I liked music
\item whether did I like music
\end{enumerate}
\end{enumerate}

\subsection*{第 4 级 · Unit 3 · 智慧学习}
\noindent\textcolor{gray}{配合 Unit 3 Smart Learning 话题的学习策略语境。}\par\smallskip
\begin{enumerate}[leftmargin=*,itemsep=4pt]
\item Our teacher told us that good learners often \_\_\_ questions.
\begin{enumerate}[label=\Alph*.,leftmargin=2em,itemsep=1pt,topsep=2pt]
\item ask
\item asked
\item will ask
\item have asked
\end{enumerate}
\item I wonder \_\_\_ I can improve my reading speed.
\begin{enumerate}[label=\Alph*.,leftmargin=2em,itemsep=1pt,topsep=2pt]
\item that
\item if
\item what
\item how
\end{enumerate}
\item Do you know \_\_\_ the library has a quiet study room?
\begin{enumerate}[label=\Alph*.,leftmargin=2em,itemsep=1pt,topsep=2pt]
\item that
\item if
\item what
\item how
\end{enumerate}
\item She asked me \_\_\_ learning strategies I used.
\begin{enumerate}[label=\Alph*.,leftmargin=2em,itemsep=1pt,topsep=2pt]
\item that
\item if
\item what
\item how
\end{enumerate}
\item I don't know \_\_\_ I should take notes in class.
\begin{enumerate}[label=\Alph*.,leftmargin=2em,itemsep=1pt,topsep=2pt]
\item that
\item whether
\item what
\item how
\end{enumerate}
\item Can you tell me \_\_\_ I can find the answer?
\begin{enumerate}[label=\Alph*.,leftmargin=2em,itemsep=1pt,topsep=2pt]
\item that
\item if
\item where
\item what
\end{enumerate}
\item He explained \_\_\_ practice makes perfect.
\begin{enumerate}[label=\Alph*.,leftmargin=2em,itemsep=1pt,topsep=2pt]
\item that
\item if
\item what
\item how
\end{enumerate}
\item I wonder \_\_\_ some students learn faster than others.
\begin{enumerate}[label=\Alph*.,leftmargin=2em,itemsep=1pt,topsep=2pt]
\item that
\item if
\item why
\item what
\end{enumerate}
\end{enumerate}

\subsection*{第 5 级 · Unit 4 · 我们的记忆}
\noindent\textcolor{gray}{配合 Unit 4 Our Memory 话题的记忆与回忆语境。}\par\smallskip
\begin{enumerate}[leftmargin=*,itemsep=4pt]
\item I remember \_\_\_ I left my keys on the desk.
\begin{enumerate}[label=\Alph*.,leftmargin=2em,itemsep=1pt,topsep=2pt]
\item that
\item if
\item what
\item how
\end{enumerate}
\item Do you know \_\_\_ I put my phone?
\begin{enumerate}[label=\Alph*.,leftmargin=2em,itemsep=1pt,topsep=2pt]
\item that
\item if
\item where
\item what
\end{enumerate}
\item She asked me \_\_\_ I could remember her name.
\begin{enumerate}[label=\Alph*.,leftmargin=2em,itemsep=1pt,topsep=2pt]
\item that
\item if
\item what
\item how
\end{enumerate}
\item I wonder \_\_\_ I forget things so easily.
\begin{enumerate}[label=\Alph*.,leftmargin=2em,itemsep=1pt,topsep=2pt]
\item that
\item if
\item why
\item what
\end{enumerate}
\item He told me \_\_\_ he had seen the film before.
\begin{enumerate}[label=\Alph*.,leftmargin=2em,itemsep=1pt,topsep=2pt]
\item that
\item if
\item what
\item how
\end{enumerate}
\item Can you tell me \_\_\_ memory works?
\begin{enumerate}[label=\Alph*.,leftmargin=2em,itemsep=1pt,topsep=2pt]
\item that
\item if
\item how
\item what
\end{enumerate}
\item I don't know \_\_\_ the story is true.
\begin{enumerate}[label=\Alph*.,leftmargin=2em,itemsep=1pt,topsep=2pt]
\item that
\item whether
\item what
\item how
\end{enumerate}
\item She asked \_\_\_ I was thinking about.
\begin{enumerate}[label=\Alph*.,leftmargin=2em,itemsep=1pt,topsep=2pt]
\item that
\item if
\item what
\item how
\end{enumerate}
\end{enumerate}

\subsection*{第 6 级 · 中考 · 综合提升}
\noindent\textcolor{gray}{语序、时态、引导词混合辨析。}\par\smallskip
\begin{enumerate}[leftmargin=*,itemsep=4pt]
\item Could you tell me \_\_\_?
\begin{enumerate}[label=\Alph*.,leftmargin=2em,itemsep=1pt,topsep=2pt]
\item where is the post office
\item where the post office is
\item where the post office was
\item where was the post office
\end{enumerate}
\item I don't know \_\_\_ he will come tonight or not.
\begin{enumerate}[label=\Alph*.,leftmargin=2em,itemsep=1pt,topsep=2pt]
\item that
\item if
\item whether
\item what
\end{enumerate}
\item She asked me \_\_\_ I had finished my homework.
\begin{enumerate}[label=\Alph*.,leftmargin=2em,itemsep=1pt,topsep=2pt]
\item that
\item if
\item what
\item how
\end{enumerate}
\item The teacher told us that the earth \_\_\_ around the sun.
\begin{enumerate}[label=\Alph*.,leftmargin=2em,itemsep=1pt,topsep=2pt]
\item moves
\item moved
\item will move
\item has moved
\end{enumerate}
\item I wonder \_\_\_ we can finish the work on time.
\begin{enumerate}[label=\Alph*.,leftmargin=2em,itemsep=1pt,topsep=2pt]
\item that
\item if
\item what
\item how
\end{enumerate}
\item He asked me \_\_\_ I was going to the concert.
\begin{enumerate}[label=\Alph*.,leftmargin=2em,itemsep=1pt,topsep=2pt]
\item that
\item whether
\item what
\item how
\end{enumerate}
\item Do you know \_\_\_ the man is?
\begin{enumerate}[label=\Alph*.,leftmargin=2em,itemsep=1pt,topsep=2pt]
\item that
\item if
\item who
\item what
\end{enumerate}
\item I don't know \_\_\_ to do next.
\begin{enumerate}[label=\Alph*.,leftmargin=2em,itemsep=1pt,topsep=2pt]
\item that
\item if
\item what
\item how
\end{enumerate}
\end{enumerate}

\newpage
\section*{答案与解析}
\small
\begin{enumerate}[leftmargin=1.6em,itemsep=2pt]
\item \textbf{L1-1 A}：think 后面的从句是陈述句，用 that 引导；口语中 that 常可省略。
\item \textbf{L1-2 B}：asked 后面的从句由一般疑问句“Could you help me?”变来，用 if / whether 引导，表示“是否”。
\item \textbf{L1-3 C}：从句缺地点状语，用 where 引导；注意后面用陈述语序 the library is。
\item \textbf{L1-4 C}：与 or not 直接连用时用 whether，不用 if；whether ... or not 表示“是否”。
\item \textbf{L1-5 A}：told me 后面的从句是陈述句，用 that 引导，that 可省略。
\item \textbf{L1-6 C}：从句缺时间状语，用 when 引导；后面是陈述语序 the meeting will start。
\item \textbf{L1-7 C}：do 缺少宾语，用 what 引导；what we should do next 表示“我们接下来该做什么”。
\item \textbf{L1-8 C}：问原因用 why；why you like the book 是陈述语序。
\item \textbf{L2-1 A}：宾语从句用陈述语序：where + 主语 he + 谓语 lives，不能把 does 提前。
\item \textbf{L2-2 B}：从句主语 the train 是第三人称单数，谓语用 leaves；语序为 when + the train + leaves。
\item \textbf{L2-3 B}：what 引导的宾语从句用陈述语序：what + he + is doing。
\item \textbf{L2-4 A}：主句 asked 是过去时，从句用 was；if 后接陈述语序 I was okay。
\item \textbf{L2-5 B}：Could/Can you tell me 是委婉请求，从句表示现在的能力，用 how I can get there；不能倒装。
\item \textbf{L2-6 B}：why 引导的宾语从句用陈述语序：why + she + was late。
\item \textbf{L2-7 A}：who 在从句中作主语，直接接谓语 wrote，不能再用 did；语序 who wrote this book。
\item \textbf{L2-8 A}：主句 asked 是过去时，从句用过去完成时 had seen；whether 后接陈述语序。
\item \textbf{L3-1 B}：主句 said 是过去时，从句表示那时的状态，用一般过去时 was。
\item \textbf{L3-2 A}：主句 asks 是一般现在时，从句按实际需要，用 am。
\item \textbf{L3-3 A}：从句是客观真理，即使主句是过去时，仍用一般现在时 travels。
\item \textbf{L3-4 B}：主句 didn't know 是过去时，从句的将来用过去将来时 would come。
\item \textbf{L3-5 B}：主句 wanted 是过去时，从句表示那时正在做，用过去进行时 was doing。
\item \textbf{L3-6 B}：主句 asked 是过去时，从句用过去完成时 had been，表示“到那时为止曾经”。
\item \textbf{L3-7 B}：直接引语用现在时，变为间接引语且主句为过去时，用 was。
\item \textbf{L3-8 B}：一般疑问句变宾语从句用 if，主句 asked 过去时，like 改为 liked，语序为陈述语序。
\item \textbf{L4-1 A}：从句表示一般的学习规律，属于普遍事实，保持一般现在时 ask。
\item \textbf{L4-2 D}：问“如何提高”，用 how 引导；how I can improve ... 是陈述语序。
\item \textbf{L4-3 B}：问“是否有”，用 if / whether；从句是陈述语序 the library has ...。
\item \textbf{L4-4 C}：问“用了哪些策略”，用 what 修饰 learning strategies；what learning strategies I used。
\item \textbf{L4-5 B}：表示“是否应该”，用 whether（也可用 if）；从句用陈述语序。
\item \textbf{L4-6 C}：问地点，用 where；where I can find the answer 是陈述语序。
\item \textbf{L4-7 A}：explained 后面的从句是陈述句，用 that 引导。
\item \textbf{L4-8 C}：问原因，用 why；why some students learn faster than others。
\item \textbf{L5-1 A}：remember 后面的从句是陈述句，用 that 引导，表示“我记得……”。
\item \textbf{L5-2 C}：问地点，用 where；where I put my phone 是陈述语序。
\item \textbf{L5-3 B}：问“是否能”，用 if / whether；从句用陈述语序 I could remember ...。
\item \textbf{L5-4 C}：问原因，用 why；why I forget things so easily。
\item \textbf{L5-5 A}：told me 后的从句是陈述句，用 that；主句过去时，从句用过去完成时 had seen。
\item \textbf{L5-6 C}：问方式，用 how；how memory works 是陈述语序。
\item \textbf{L5-7 B}：表示“是否”，用 whether / if；从句 the story is true 是陈述语序。
\item \textbf{L5-8 C}：about 缺宾语，用 what；what I was thinking about。
\item \textbf{L6-1 B}：宾语从句用陈述语序：where + the post office + is；不能倒装。
\item \textbf{L6-2 C}：与 or not 连用用 whether；whether he will come tonight or not。
\item \textbf{L6-3 B}：问“是否”，用 if / whether；主句过去时，从句用过去完成时 had finished。
\item \textbf{L6-4 A}：客观真理用一般现在时，即使主句是 told。
\item \textbf{L6-5 B}：表示“是否能”，用 if / whether；if we can finish ... 是陈述语序。
\item \textbf{L6-6 B}：问“是否”，用 whether / if；从句用陈述语序 I was going ...。
\item \textbf{L6-7 C}：问“那个人是谁”，用 who；who the man is 是陈述语序。
\item \textbf{L6-8 C}：do 缺宾语，用 what；what to do next 是“疑问词 + 不定式”结构。
\end{enumerate}

\end{document}
